\documentclass[11pt,a4paper]{article}
\usepackage[margin=1.05in]{geometry}
\usepackage{amsmath,amssymb,amsthm}
\usepackage{microtype}
\usepackage[hidelinks]{hyperref}
\makeatletter
\g@addto@macro{\UrlBreaks}{\do\_}\makeatother
\hypersetup{pdftitle={When gold disconnects},pdfauthor={Jack Pickett}}

\newtheorem{lemma}{Lemma}
\newtheorem{corollary}{Corollary}
\newtheorem{remark}{Remark}
\newtheorem{theorem}{Theorem}

\title{When gold disconnects}
\author{Jack Pickett \textperiodcentered\ 19th September 2026}
\date{}

\begin{document}
\maketitle
\thispagestyle{empty}

\begin{abstract}
In the hire graph of the 3-free door, every odd prime other than \(3\) meets \(5\) in the infinite undirected gold graph.
Finite windows can still be disconnected: each new \(2\)-power-door prime is gold-isolated at its first-owner window.
Downward gold chains terminate only at \(2\)-power doors, so gold is disconnected for infinitely many finite \(X\) if and only if there are infinitely many Mersenne or Fermat primes.
The companion note settles the infinite component of \(5\); this note records the equivalence for the leftover finite-window question, together with thin computational certificates for the \(M_{31}\) corridor.
\end{abstract}

\section{Setup}

We use the notation of the hire-graph note~\cite{Hire}: the 3-free door \(m_0\), the hire set \(S_X\), directed gold arcs \(q\mid m_0(p)\), and the undirected gold graph on \(S_X\setminus\{2\}\).
A prime \(M\neq 3\) is a \emph{\(2\)-power door} if \(m_0(M)=2^k\) for some \(k\ge 1\).
Every Mersenne prime \(M_p=2^p-1\) with \(p\) odd, and every Fermat prime greater than \(3\), is of this form.
The first owner \(r_0\) of \(M\) is the least prime \(r\) with \(M\mid m_0(r)\); by the door bound, \(r_0\ge 2M-1\).

\section{First-owner windows}

\begin{lemma}
\label{lem:isolated}
Let \(M>5\) be a \(2\)-power-door prime and let \(r_0\) be its first owner.
At the window \(X=r_0\), the prime \(M\) is an isolated vertex of undirected gold.
In particular gold is disconnected.
\end{lemma}

\begin{proof}
Since \(M\mid m_0(r_0)\) and \(r_0\le X\), one has \(M\in S_X\).
The first owner itself is not hired at this window: nothing smaller than \(r_0\) can have hired it, so \(r_0\notin S_X\).
As a \(2\)-power door, \(M\) has gold out-degree zero.
No other owner of \(M\) exists below \(r_0\), so no incoming gold arc is live either.
Thus \(M\) meets no other odd hired prime by a gold edge.
Since \(M>5\) and \(r_0\ge 2M-1>11\), we also have \(5\in S_{r_0}\) (for example \(5\mid m_0(11)=10\)), so \(M\) is an isolated vertex distinct from another vertex of the gold graph.
\end{proof}

As of September 2026, every known Mersenne prime from \(M_{31}\) onward has first owner larger than the connecting window \(X^*=92\,274\,421\) of the hire-graph note, so Lemma~\ref{lem:isolated} already supplies many concrete disconnected windows.
That is a finite list, not an infinitude theorem.

\section{Downward chains}

\begin{lemma}[Door closure in a window]
\label{lem:closed}
If \(q\in S_X\setminus\{2\}\) and \(r\) is an odd prime dividing \(m_0(q)\), then \(r\in S_X\).
\end{lemma}

\begin{proof}
Every odd hired prime satisfies \(q\le(X+1)/2\le X\), so \(q\) itself is an owner in the window.
Any odd prime factor of \(m_0(q)\) is therefore hired into \(S_X\).
\end{proof}

\begin{lemma}
\label{lem:sinks}
Let \(S\) be a finite nonempty set of odd primes other than \(3\), closed under taking odd prime factors of \(m_0\).
Then \(S\) contains a \(2\)-power door.
\end{lemma}

\begin{proof}
Take a least element \(p\in S\).
If \(m_0(p)=2^k\), we are done.
Otherwise some odd prime \(q\neq 3\) divides \(m_0(p)\).
Closure puts \(q\in S\), and the size bound \(q\le m_0(p)/2<(p+1)/2\) forces \(q<p\), contradicting minimality.
\end{proof}

Gold arcs on odd leaves are strictly decreasing in the prime label, so every finite descending gold chain ends at a sink, and Lemma~\ref{lem:sinks} says every such sink in a door-closed set is a \(2\)-power door.

\section{The equivalence}

\begin{theorem}
\label{thm:equiv}
The following are equivalent.
\begin{enumerate}
\item There are infinitely many finite windows \(X\) on which undirected gold is disconnected.
\item There are infinitely many \(2\)-power-door primes.
\item There are infinitely many Mersenne primes, or infinitely many Fermat primes.
\end{enumerate}
\end{theorem}

\begin{proof}
Lemma~\ref{lem:isolated} gives \((2)\Rightarrow(1)\): each sufficiently large \(2\)-power door contributes at least its first-owner window.
For \((1)\Rightarrow(2)\), suppose only finitely many \(2\)-power doors exist.
After the last of them has been absorbed into the infinite component of \(5\) (which the hire-graph note supplies for every odd prime other than \(3\)), every later hired vertex has a downward gold path onto an already-absorbed sink.
That path is live in the window by Lemma~\ref{lem:closed}.
Hence later windows stay connected.
The equivalence of \((2)\) and \((3)\) is the classification of \(2\)-power doors as Mersenne or Fermat primes greater than \(3\).
\end{proof}

Fermat primes are widely believed finite (five known; the last is \(65537\)).
Infinitude of Mersenne primes is open.
Thus Theorem~\ref{thm:equiv} does not settle item~(1) unconditionally; it identifies the obstruction.

\begin{remark}[Heuristic]
The Lenstra--Pomerance--Wagstaff heuristic predicts about \(e^\gamma\log_2 x\) Mersenne primes with exponent at most \(x\), hence infinitely many, and therefore infinitely many disconnected windows.
(An additional, separate heuristic is that each large Mersenne prime opens a satellite interval whose length is governed by owner and hire delays in residue classes modulo a multiple of \(M_p\); that is not a consequence of LPW alone.)
That is a labelled heuristic, not a proof of~(1).
\end{remark}

\section{The \(M_{31}\) integers}

Let \(M_{31}=2^{31}-1\).
Thin CRT searches (exhaustive over the two admissible classes, not a gold-graph rebuild) yield the following computational certificates:
\[
r_0=98\,784\,247\,763,\qquad
v=300\,647\,710\,579,
\]
where \(r_0\) is certified as the least prime owner of \(M_{31}\) in either class modulo \(3M_{31}\), and \(v\) is certified as the least prime in either class modulo \(15M_{31}\) with \(5M_{31}\mid m_0(v)\) (class A; \(m_0(v)=2^2\cdot 5\cdot 7\cdot M_{31}\)).
The ratio \(v/r_0\approx 3.04\).
The floor candidate \(2M_{31}-1\) is composite (\(9241\cdot 464773\)); the certificate for \(r_0\) is that no earlier candidate in those two classes is prime.

At \(X=v\) the integer \(v\) is an owner, so \(5\) and \(M_{31}\) divide its door, but \(v\notin S_X\), and the gold arcs \(v\to 5\) and \(v\to M_{31}\) are not live until \(v\) itself is hired.
That hiring time is at least \(2v-1\approx 6.01\cdot 10^{11}\).
What \(v\) records is a bridge owner, not yet a bridge vertex.
A mixed owner hired earlier than \(v\) could still connect \(M_{31}\) sooner; ``first direct \(5\)--\(M_{31}\) bridge'' is not automatically the first connecting vertex.
A full component check at these scales remains past home computing.

\section*{Lean}

Lemmas~\ref{lem:isolated}--\ref{lem:sinks} and Theorem~\ref{thm:equiv} are checked in Lean~4
(\path{Hire/GoldDisconnects.lean}; declaration names on the companion site).
Door closure (Lemma~\ref{lem:closed}) is the next Lean target.
The \(M_{31}\) integers above are thin computational certificates, not Lean theorems.

\begin{thebibliography}{9}

\bibitem{Hire}
J.~Pickett,
\emph{The hire graph of the 3-free door},
note, 19th September 2026.

\end{thebibliography}

\end{document}
